\section{Introduction}

Smartphone recordings of cough, breathing, and speech offer a convenient way
to study respiratory health outside a clinic. Public datasets and early model
studies therefore motivated audio-based COVID-19 screening as a possible
low-cost triage aid \cite{xia2021covid,bhattacharya2023coswara,orlandic2021coughvid}.
The central methodological difficulty is that these datasets were assembled
during a rapidly changing pandemic through heterogeneous recruitment,
self-report, devices, locations, and testing practices. A classifier can
separate labels in one sample without learning a stable disease-related audio
pattern that transfers to another period or collection system.

Several strands of prior work illustrate the tension. Architecture-led studies
reported strong within-dataset discrimination using transfer learning,
multimodal score fusion, and hierarchical spectrogram transformers
\cite{pahar2022transfer,chetupalli2023multimodal,aytekin2024hst}. In contrast, realistic
participant-independent evaluation showed that biased split construction can
raise AUROC from approximately 0.71 to 0.90 \cite{han2022realistic}. Coppock
\emph{et al.} found that an audio model's AUROC decreased from 0.846 to 0.619
after matching measured confounders, and that symptom models were strong
comparators \cite{coppock2024audio}. Temporal drift and cross-dataset studies
have since reinforced that within-dataset performance and deployment-time
transportability are different estimands \cite{ganitidis2025drift,islam2026robust}.

This study asks a narrower question than whether one more classifier can raise
an internal benchmark. It asks which conclusions survive when models developed
from the same source data are examined under different sources of shift:
participant allocation, calendar time, label prevalence, recording protocol,
and dataset origin. We use Coswara because it contains multiple sound types per
participant, and COUGHVID as a separately collected cough-only target. We also
test whether non-audio variables predict the label, whether audio adds value to
those variables on aligned participants, and whether calibration or target
threshold adjustment can repair weak external discrimination.

The contributions are:

\begin{enumerate}
    \item a participant-audited multimodal pipeline spanning conventional
    acoustic functionals and learned audio representations;
    \item protocol-specific internal, calendar-aware, chronological, and
    modality-matched external evaluations with uncertainty estimates;
    \item negative controls, metadata prediction, feature-stability, and
    domain-overlap analyses that distinguish observed associations from
    stronger causal interpretations; and
    \item base-rate-aware calibration, fixed-sensitivity, and decision analyses
    that separate ranking performance from operational usefulness.
\end{enumerate}

The study is not a clinical validation and does not claim that COVID-19 leaves
no acoustic signature. Its purpose is to determine how much evidence the
available retrospective benchmarks provide for a reliable screening model.
